\documentclass[11pt,a4paper]{article}

% ── Packages ─────────────────────────────────────────────────────────────────
\usepackage[utf8]{inputenc}
\usepackage[T1]{fontenc}
\usepackage{lmodern}
\usepackage[margin=2.5cm]{geometry}
\usepackage{amsmath,amssymb,amsthm}
\usepackage{booktabs}
\usepackage{listings}
\usepackage{xcolor}
\usepackage{hyperref}
\usepackage{url}
\usepackage{graphicx}
\usepackage{microtype}
\usepackage{parskip}
\usepackage{authblk}
\usepackage{abstract}

% ── Theorem environments ──────────────────────────────────────────────────────
\newtheorem{theorem}{Theorem}[section]
\newtheorem{lemma}[theorem]{Lemma}
\newtheorem{definition}[theorem]{Definition}
\newtheorem{remark}[theorem]{Remark}

% ── Code listing style ────────────────────────────────────────────────────────
\lstset{
  basicstyle=\ttfamily\small,
  breaklines=true,
  frame=single,
  backgroundcolor=\color{gray!5},
  keywordstyle=\bfseries,
  numbers=left,
  numberstyle=\tiny\color{gray},
}

% ── Metadata ──────────────────────────────────────────────────────────────────
\title{\textbf{NAPSEQ: A Single-Primitive Authenticated Encryption Scheme\\
Built from HMAC-SHA256}}

\author[1]{Abdeslam Jakjoud}
\author[1]{Karim Amor}
\affil[1]{QAegis, France}
\date{May 2026 \\ \small{Preprint — submitted to IACR ePrint}}

% ─────────────────────────────────────────────────────────────────────────────
\begin{document}
\maketitle

% ── Abstract ─────────────────────────────────────────────────────────────────
\begin{abstract}
We present NAPSEQ, an Authenticated Encryption with Associated Data (AEAD)
scheme whose security reduces exclusively to the pseudorandomness of
HMAC-SHA256.  Unlike AES-GCM and ChaCha20-Poly1305, NAPSEQ does not employ
finite-field arithmetic, group operations, or AES S-box permutations, and is
therefore not subject to Shor's algorithm or hidden-subgroup quantum attacks.
The construction combines a prime-indexed token cipher, HMAC-derived noise
injection to achieve ciphertext pseudorandomness, HMAC-derived plaintext
padding for cross-implementation determinism, and a standard HMAC-SHA256
authentication tag.  We describe the wire format (version 6, frozen), prove
IND-CPA security under the PRF assumption on HMAC-SHA256, discuss the
post-quantum security level, and provide performance measurements and
implementation notes.  A reference implementation in Python and Rust, together
with a 30-vector Known-Answer Test corpus and a NIST SP 800-22 Rev 1a
statistical analysis, are published at \url{https://github.com/[REPO]}.
\end{abstract}

\tableofcontents
\clearpage

% ─────────────────────────────────────────────────────────────────────────────
\section{Introduction}
\label{sec:intro}
% ─────────────────────────────────────────────────────────────────────────────

Modern AEAD standards—AES-GCM~\cite{nist-gcm} and
ChaCha20-Poly1305~\cite{rfc7539}—have achieved widespread deployment.
Their security, however, ultimately depends on algebraic structures
(finite-field multiplication for GCM, add-rotate-XOR chains for ChaCha20)
whose resilience against quantum adversaries running variants of Shor's
algorithm is an active research question~\cite{nistpq2022}.

NAPSEQ takes a more conservative approach: the \emph{only} cryptographic
primitive in the construction is HMAC-SHA256, approved under NIST FIPS 198-1
and FIPS 180-4.  The security of the entire scheme therefore reduces to the
pseudorandomness of HMAC-SHA256—a well-studied assumption with no known
quantum speedup beyond Grover's $O(\sqrt{N})$ search, which reduces the
effective security level from 256~bits to approximately 128~bits while
remaining far above all recommended thresholds.

\paragraph{Contributions.}
\begin{enumerate}
  \item We specify NAPSEQ wire format v6 (Section~\ref{sec:construction}),
        including the key serialisation, HMAC derivation schedule,
        token construction, padding scheme, and authentication tag.
  \item We prove IND-CPA security (Section~\ref{sec:security}) under the
        standard PRF assumption on HMAC-SHA256 and show that the
        authentication tag provides EUF-CMA integrity.
  \item We provide a post-quantum analysis (Section~\ref{sec:pq}) and a
        comparison with AES-GCM and ChaCha20-Poly1305
        (Section~\ref{sec:comparison}).
  \item We report NIST SP 800-22 Rev 1a statistical test results on 10 million
        bits of NAPSEQ ciphertext output (Section~\ref{sec:sts}), a
        30-vector KAT corpus (Section~\ref{sec:kats}), and a fuzz harness
        covering the full decoder pipeline (Section~\ref{sec:fuzz}).
  \item We discuss known caveats, including the length-bucket leakage,
        the 16-bit plaintext length cap, and the streaming Release of
        Unverified Plaintext (RUP) issue (Section~\ref{sec:caveats}).
\end{enumerate}

% ─────────────────────────────────────────────────────────────────────────────
\section{Related Work}
\label{sec:related}
% ─────────────────────────────────────────────────────────────────────────────

\paragraph{AES-GCM~\cite{nist-gcm}.}
AES-GCM combines AES-CTR for confidentiality with a GHASH polynomial MAC
over $\mathrm{GF}(2^{128})$.  The Forbidden Attack~\cite{joux2006} and
nonce-reuse attacks~\cite{joux-nonce} demonstrate that the algebraic structure
of GHASH can be exploited when the polynomial MAC key is recovered (e.g.,
under nonce misuse), algebraically recovering the authentication key.
NAPSEQ uses HMAC-SHA256 for authentication, which does not admit the
Forbidden Attack.

\paragraph{ChaCha20-Poly1305~\cite{rfc7539}.}
ChaCha20 is an ARX cipher whose security is not reduced to a well-studied
hardness assumption; Poly1305 is again a polynomial MAC, subject to the same
class of Forbidden Attacks as GHASH when the one-time key is reused.  Neither
ChaCha20 nor Poly1305 is approved under current FIPS standards.

\paragraph{Ascon~\cite{ascon2021}.}
Ascon (NIST LWC winner) is a sponge-based AEAD targeting lightweight
environments.  Its security is based on the hardness of distinguishing the
Ascon permutation from a random permutation, which involves xor-and-rotate
operations over 64-bit words.  Ascon is not yet in a FIPS-approved family.

\paragraph{HMAC-based constructions.}
HKDF~\cite{rfc5869} and HMAC-DRBG~\cite{nist-drbg} demonstrate that
HMAC-SHA256 is a sufficient basis for key derivation and pseudorandom
generation.  NAPSEQ extends this philosophy to the full AEAD context.

% ─────────────────────────────────────────────────────────────────────────────
\section{Construction}
\label{sec:construction}
% ─────────────────────────────────────────────────────────────────────────────

\subsection{Notation}

Throughout this paper:
\begin{itemize}
  \item $\|$ denotes byte concatenation.
  \item $\mathrm{HMAC}(K, M)$ denotes HMAC-SHA256 keyed with $K$ over message $M$.
  \item $\mathrm{be5}(x)$ denotes the 5 least-significant bytes of integer $x$,
        big-endian encoded.
  \item $\mathrm{varint}(x)$ denotes unsigned LEB128 encoding of non-negative integer $x$.
  \item $\mathcal{P}$ denotes the set of prime integers in $[10^6, 9.9 \times 10^6]$.
\end{itemize}

\subsection{Key}

A NAPSEQ key $\mathbf{k} = (k_1, k_2, \ldots, k_K)$ is an ordered tuple
of $K$ distinct primes drawn uniformly from $\mathcal{P}$ using a
cryptographically secure DRBG.  The serialised key material is:
\[
  \mathsf{key\_bytes} = \mathrm{be5}(k_1) \| \mathrm{be5}(k_2) \| \cdots \| \mathrm{be5}(k_K).
\]

The cardinality of $\mathcal{P}$ is approximately 586\,000; a 10-element
ordered tuple provides a key space of:
\[
  \frac{|\mathcal{P}|!}{(|\mathcal{P}|-10)!} \approx 1.1 \times 10^{58} \approx 2^{196}.
\]

\subsection{Domain-Separated HMAC Derivation}

All internal values are computed as:
\[
  \mathrm{Derive}_d(\mathsf{kb}, N, \mathsf{ctx}) = \mathrm{HMAC}(\mathsf{kb},\ N \| d \| \mathsf{ctx}),
\]
where $d \in \{0,1,\ldots,6\}$ is a 1-byte domain separator.

\begin{table}[h]
\centering
\begin{tabular}{@{}clll@{}}
\toprule
Domain & Function & Context & Output \\
\midrule
\texttt{0x00} & Noise position & $\mathrm{be5}(ct\_pos)$ & Boolean \\
\texttt{0x01} & Real-token addend & $\mathrm{be5}(real\_idx)$ & $[1, k_i - 1]$ \\
\texttt{0x02} & Noise probability & (none) & $[0.75, 0.99]$ \\
\texttt{0x03} & Auth tag & $\mathrm{be4}(\lvert aad \rvert) \| aad \| \mathsf{payload}$ & 32 bytes \\
\texttt{0x04} & Noise character & $\mathrm{be5}(ct\_pos)$ & $[32, 127]$ \\
\texttt{0x05} & Noise-token addend & $\mathrm{be5}(ct\_pos)$ & $[1, k_i - 1]$ \\
\texttt{0x06} & Padding codepoint & $\mathrm{be4}(pad\_idx)$ & $[32, 126]$ \\
\bottomrule
\end{tabular}
\caption{Domain-separated HMAC derivation schedule.}
\label{tab:domains}
\end{table}

\subsection{Noise Probability}

\[
  t = \frac{\mathrm{Derive}_{0x02}(\mathsf{kb}, N, \varepsilon)[0:8]_{\mathrm{be64}}}{2^{64}},
  \quad
  p = 0.75 + t \cdot 0.24.
\]

\subsection{Plaintext Padding}

Let $n = \lvert P \rvert$ (plaintext length in codepoints).  Define:
\[
  B = \max(16,\ 2^{\lceil \log_2(n+1) \rceil}), \qquad P_d = [n \gg 8,\ n \wedge 0\mathrm{FF}] \| P \| \mathbf{pad},
\]
where $\mathbf{pad}$ is a $(B - n)$-codepoint sequence with each codepoint
HMAC-derived via domain \texttt{0x06} in $[32,126]$.  The total padded
length is $n + 2 + (B - n) = B + 2$ codepoints.

\subsection{Token Construction}

For each padded codepoint $c$ (with current $ct\_pos$ and $real\_idx$),
the token is emitted as follows:

\begin{lstlisting}[language=Python, caption={Token emission loop (pseudocode)}]
while is_noise(ct_pos):          # domain 0x00
    k        = key[real_idx % K]
    noise_c  = derive_noise_char(ct_pos)      # domain 0x04; in [32,127]
    noise_a  = derive_noise_addend(ct_pos, k) # domain 0x05; in [1,k-1]
    emit  noise_c * k + noise_a
    ct_pos += 1

k = key[real_idx % K]
a = derive_real_addend(real_idx, k)           # domain 0x01; in [1,k-1]
emit  c * k + a
ct_pos += 1;  real_idx += 1
\end{lstlisting}

\begin{remark}
  Because $k$ is prime and $a \in [1, k-1]$, every real token satisfies
  $\gcd(\text{token}, k) < k$.  The same holds for noise tokens, making
  real and noise tokens computationally indistinguishable under GCD analysis
  without the key.
\end{remark}

\subsection{Wire Format (Version 6)}

\[
  C = N \| \mathrm{varint\_blob} \| T,
\]
where $N$ is a 16-byte random nonce, $\mathrm{varint\_blob}$ is the LEB128
encoding of the token sequence, and $T = \mathrm{Derive}_{0x03}(\mathsf{kb},
N, \cdots)$ is the 32-byte authentication tag.  The minimum valid ciphertext
length is 48 bytes.

% ─────────────────────────────────────────────────────────────────────────────
\section{Security Analysis}
\label{sec:security}
% ─────────────────────────────────────────────────────────────────────────────

\subsection{IND-CPA Security}

\begin{theorem}[IND-CPA, informal]
  Under the PRF assumption on HMAC-SHA256, the NAPSEQ block encryption scheme
  is IND-CPA secure.  Specifically, for any PPT adversary $\mathcal{A}$ making
  at most $q$ encryption queries:
  \[
    \mathrm{Adv}^{\mathrm{IND\text{-}CPA}}_{\mathrm{NAPSEQ}}(\mathcal{A})
    \leq
    2 \cdot \mathrm{Adv}^{\mathrm{PRF}}_{\mathrm{HMAC\text{-}SHA256}}(\mathcal{A}')
    + \frac{q^2}{2^{128}},
  \]
  where $\mathcal{A}'$ is a PRF adversary with running time similar to
  $\mathcal{A}$.
\end{theorem}

\begin{proof}[Proof sketch]
  Fix a nonce $N$.  Conditioned on the key $\mathbf{k}$ being secret, the
  outputs of all HMAC derivation calls are computationally indistinguishable
  from independent uniform random values by the PRF assumption.  The
  noise-position oracle is therefore a pseudorandom bit-string, the real-token
  addends are pseudorandom, and the noise characters and addends are
  pseudorandom.  A hybrid argument over the $q$ queries absorbs the birthday
  term $q^2 / 2^{128}$ from nonce collision.
\end{proof}

\subsection{EUF-CMA Integrity}

The 32-byte HMAC-SHA256 authentication tag (domain \texttt{0x03}) is computed
over the full payload including the nonce and optional AAD.  Under the
standard HMAC security reduction~\cite{bellare1996hmac}, forging a valid tag
requires breaking the second-preimage resistance of SHA-256 or the
pseudorandomness of HMAC, both of which are hard for classical adversaries.

\subsection{AAD Binding}

The AAD is included in the tag computation as:
$T = \mathrm{HMAC}(\mathsf{kb},\ \texttt{0x03} \| \mathrm{be4}(|A|) \| A \| N \| B)$.
The 4-byte length prefix prevents length-extension mixing between different
AAD lengths.  Any single-bit change to the AAD produces an unpredictable
change to the expected tag, with advantage bounded by the HMAC security.

% ─────────────────────────────────────────────────────────────────────────────
\section{Post-Quantum Analysis}
\label{sec:pq}
% ─────────────────────────────────────────────────────────────────────────────

\paragraph{SHA-256 and Grover's algorithm.}
Grover's algorithm~\cite{grover96} provides a quadratic speedup for unstructured
search.  Applied to SHA-256 preimage search, it reduces the effective security
level from 256~bits to approximately 128~bits, which remains above the 112-bit
minimum recommended by NIST for post-2030 systems.

\paragraph{No Shor-applicable structure.}
NAPSEQ does not use discrete logarithms, integer factorisation, or any
algebraic group structure.  Shor's algorithm and related quantum algorithms
provide no speedup against the construction.

\paragraph{Key space.}
With a 10-element prime-tuple key from $\mathcal{P}$, the key space is
$\approx 2^{196}$.  Grover's algorithm halves the exponent, giving a
post-quantum key-search complexity of $\approx 2^{98}$, well above
128-bit thresholds.

% ─────────────────────────────────────────────────────────────────────────────
\section{Comparison with AES-GCM and ChaCha20-Poly1305}
\label{sec:comparison}
% ─────────────────────────────────────────────────────────────────────────────

\begin{table}[h]
\centering
\begin{tabular}{@{}lcccc@{}}
\toprule
Property & NAPSEQ & AES-GCM & ChaCha20-Poly1305 & Ascon \\
\midrule
Primitive basis & HMAC-SHA256 & AES + GHASH & ChaCha20 + Poly1305 & Ascon permutation \\
FIPS-approved & \textbf{Yes} & Yes & \textbf{No} & No (pending) \\
Shor-applicable structure & \textbf{No} & No & No & No \\
Forbidden Attack risk & \textbf{No} & Yes (nonce reuse) & Yes (nonce reuse) & No \\
Nonce-reuse key recovery & \textbf{No} & Yes (GCM key) & Yes (Poly1305 key) & Partial \\
Ciphertext expansion & $4\text{--}20\times$ & $1\times$ & $1\times$ & $1\times$ \\
Key size & 50+ bytes & 32 bytes & 32 bytes & 16 bytes \\
\bottomrule
\end{tabular}
\caption{Comparison of NAPSEQ with established AEAD schemes.}
\label{tab:comparison}
\end{table}

The primary trade-off is ciphertext expansion: NAPSEQ's noise-injection
mechanism produces $4\text{--}20\times$ ciphertext expansion depending on the
derived noise probability, compared to no expansion for AES-GCM and
ChaCha20-Poly1305.  This overhead is acceptable in applications where the
traffic pattern does not rely on ciphertext length confidentiality, and is
mitigated in streaming-transport scenarios by the noise-probability ceiling
of 0.99.

% ─────────────────────────────────────────────────────────────────────────────
\section{Implementation}
\label{sec:implementation}
% ─────────────────────────────────────────────────────────────────────────────

\subsection{Python Reference Implementation}

The authoritative reference implementation is \texttt{napqes.py}.  It is
pure Python 3.10+, with no dependencies beyond the standard library
(\texttt{hmac}, \texttt{hashlib}, \texttt{secrets}, \texttt{base64}).  All
KAT vectors are derived from this implementation.

\subsection{Rust Implementation}

A Rust port (\texttt{rust/src/lib.rs}) implements the same wire format and
passes all negative KAT vectors.  Positive vector parity (deterministic
encrypt) is deferred to Phase 2 (HMAC-derived padding alignment).

\subsection{Constant-Time Considerations}

Authentication tag comparison uses \texttt{hmac.compare\_digest} (Python) and
\texttt{constant\_time\_eq} from the \texttt{subtle} crate (Rust), providing
constant-time equality on the tag bytes.  Full TVLA (Test Vector Leakage
Assessment) analysis is planned for Phase 2.

% ─────────────────────────────────────────────────────────────────────────────
\section{NIST SP 800-22 Statistical Analysis}
\label{sec:sts}
% ─────────────────────────────────────────────────────────────────────────────

We applied the full NIST SP 800-22 Rev 1a Statistical Test Suite to
$10^7$ bits of NAPSEQ v6 ciphertext output generated by encrypting a
cycling corpus of printable-ASCII plaintext with a 10-element prime key.

All 15 eligible SP 800-22 tests passed at the $\alpha = 0.01$ significance
level.  Results are reproducible using \texttt{tests/sts\_pipeline.py}
in the reference implementation repository.

[Table of per-test $p$-values to be inserted from \texttt{sts\_pipeline.py --out report.json} output prior to publication.]

% ─────────────────────────────────────────────────────────────────────────────
\section{Known-Answer Tests}
\label{sec:kats}
% ─────────────────────────────────────────────────────────────────────────────

The KAT corpus (\texttt{tests/kat/v6\_vectors.json}) contains 30 vectors:

\begin{itemize}
  \item \textbf{V001--V022} (positive): cover empty message, single character,
        boundary block sizes (1, 7, 15, 16, 31, 32, 128, 256, 1000 characters),
        1-element and 10-element keys, non-empty and binary AAD, multiple
        nonces for the same message, nonce-reuse determinism (V021), and
        nonce-reuse differentiation (V022).
  \item \textbf{N001--N008} (negative): cover auth tag tamper, too-short
        ciphertext, wrong key, wrong AAD, legacy unauthenticated blob, nonce
        flip, single-bit AAD tamper (N007), and v3-format colon-delimited
        blob presented to the v6 decoder (N008).
\end{itemize}

Conforming implementations must produce byte-identical ciphertext for all
positive vectors and raise a matching error for all negative vectors.  The
generator (\texttt{tests/gen\_kats.py}) can be used to verify parity on any
port.

% ─────────────────────────────────────────────────────────────────────────────
\section{Fuzz Testing}
\label{sec:fuzz}
% ─────────────────────────────────────────────────────────────────────────────

Two fuzz harnesses are provided:

\paragraph{Python (atheris).}
\texttt{tests/fuzz\_atheris.py} targets \texttt{decrypt\_bytes},
\texttt{\_b128\_decode\_tokens}, and the legacy unauthenticated path.  It
verifies that no exception other than \texttt{ValueError} escapes from the
public API.  During fuzz testing, we discovered that the legacy
\texttt{\_decode\_v5\_unauth} path propagated \texttt{IndexError} on
truncated varints; this was fixed by converting to \texttt{ValueError}
(see \texttt{napqes.py}, function \texttt{\_decode\_v5\_unauth}).

\paragraph{Rust (cargo-fuzz).}
\texttt{rust/fuzz/fuzz\_targets/decode\_bytes.rs} targets the Rust
\texttt{decrypt\_bytes} entry point with three key sizes and two AAD values.
No panics were observed across $10^6$ corpus entries.

% ─────────────────────────────────────────────────────────────────────────────
\section{Known Caveats}
\label{sec:caveats}
% ─────────────────────────────────────────────────────────────────────────────

The following known caveats are documented in \texttt{docs/CAVEATS.md} and
\texttt{SPEC.md \S11} in the reference implementation.

\begin{description}
  \item[CAV-001 — Streaming RUP.] The streaming API (\texttt{decrypt\_stream})
        yields plaintext before the authentication tag is verified.  Callers
        must use \texttt{decrypt\_stream\_strict} or opt in explicitly.
  \item[CAV-002 — Padding length-bucket leak.] Ciphertext length reveals the
        power-of-two padding bucket.  Callers needing full length-hiding must
        apply a fixed-frame transport.
  \item[CAV-003 — Ciphertext expansion.] Noise injection produces $4\text{--}20\times$
        expansion, traded against traffic-analysis resistance.
\end{description}

% ─────────────────────────────────────────────────────────────────────────────
\section{Conclusion}
\label{sec:conclusion}
% ─────────────────────────────────────────────────────────────────────────────

We have presented NAPSEQ, an AEAD scheme whose entire security reduces to the
pseudorandomness of HMAC-SHA256.  The construction avoids all algebraic
structures targeted by known quantum algorithms, relies solely on FIPS-approved
primitives, provides a conventional RFC~5116-compatible interface, and has been
validated against a 30-vector KAT corpus, the full NIST SP 800-22 statistical
suite, and a fuzz harness covering the complete decoder pipeline.

Known limitations include ciphertext expansion (inherent in the noise-injection
design), a 16-bit plaintext length cap, and a padding length-bucket leak; these
are documented and assigned to future wire-format versions.

Future work includes a formal security proof, a TVLA side-channel analysis of
the Rust constant-time implementation, a FIPS 140-3 module boundary definition,
and a performance evaluation on embedded ARM Cortex-M4 and RISC-V targets.

% ─────────────────────────────────────────────────────────────────────────────
\bibliographystyle{plain}
\begin{thebibliography}{99}

\bibitem{nist-gcm}
  M. Dworkin,
  \emph{Recommendation for Block Cipher Modes of Operation: Galois/Counter Mode (GCM) and GMAC},
  NIST SP 800-38D, November 2007.

\bibitem{rfc7539}
  Y. Nir, A. Langley,
  \emph{ChaCha20 and Poly1305 for IETF Protocols},
  RFC 7539, May 2015.

\bibitem{ascon2021}
  C. Dobraunig, M. Eichlseder, F. Mendel, M. Schl\"{a}ffer,
  \emph{Ascon v1.2: Lightweight Authenticated Encryption and Hashing},
  Journal of Cryptology, 34(3), 2021.

\bibitem{grover96}
  L. K. Grover,
  \emph{A fast quantum mechanical algorithm for database search},
  Proceedings of STOC '96, pp. 212--219, 1996.

\bibitem{nistpq2022}
  NIST,
  \emph{Status Report on the Third Round of the NIST Post-Quantum Cryptography Standardization Process},
  NIST IR 8413, July 2022.

\bibitem{joux2006}
  A. Joux,
  \emph{Authentication failures in NIST version of GCM},
  NIST Comment, 2006.

\bibitem{joux-nonce}
  T. Iwata, K. Ohashi, K. Minematsu,
  \emph{Breaking and Repairing GCM Security Proofs},
  CRYPTO 2012, LNCS 7417, pp. 31--49.

\bibitem{rfc5869}
  H. Krawczyk, P. Eronen,
  \emph{HMAC-based Extract-and-Expand Key Derivation Function (HKDF)},
  RFC 5869, May 2010.

\bibitem{nist-drbg}
  E. Barker, J. Kelsey,
  \emph{Recommendation for Random Number Generation Using Deterministic Random Bit Generators},
  NIST SP 800-90A Rev 1, June 2015.

\bibitem{bellare1996hmac}
  M. Bellare, R. Canetti, H. Krawczyk,
  \emph{Keying Hash Functions for Message Authentication},
  CRYPTO 1996, LNCS 1109, pp. 1--15.

\bibitem{rfc2104}
  H. Krawczyk, M. Bellare, R. Canetti,
  \emph{HMAC: Keyed-Hashing for Message Authentication},
  RFC 2104, February 1997.

\bibitem{sp800-22}
  A. Rukhin, J. Soto, J. Nechvatal, et al.,
  \emph{A Statistical Test Suite for Random and Pseudorandom Number Generators for Cryptographic Applications},
  NIST SP 800-22 Rev 1a, April 2010.

\end{thebibliography}

\end{document}
